% ============================================================
% Appendix A. Reproducibility details  (article-class version)
% Included after \bibliography{references}
% ============================================================
\appendix
\section*{Appendix A. Reproducibility details}
\addcontentsline{toc}{section}{Appendix A. Reproducibility details}

The reproducibility details below are taken directly from the implementation and saved experimental outputs. Relevant source files include \texttt{code.py} (data-generating processes, base learners, and metrics), \texttt{task\_common.py} (data loading, cross-fitted encoder, conformal threshold, and bootstrap procedure),
\texttt{task1\_cate\_benchmark.py}, \texttt{task2\_conformal.py},
\texttt{task3\_llm\_verifier.py}, \texttt{llm\_narrative\_agent.py},
\texttt{freetext\_scoring.py} and \texttt{sensitivity\_analysis.py}.

\subsection*{A.1 Cohort and design matrices}

Source table: \texttt{tcga\_ov\_master\_ml.csv} (TCGA-OV). Fields used:
\texttt{age\_at\_diagnosis}, \texttt{figo\_stage}, \texttt{race},
\texttt{ethnicity} (primary, four fields).
The GDC missing marker \texttt{'--} is mapped to
\texttt{NaN} and those rows are dropped; \texttt{"not reported"} is retained as
a legitimate category level. Age is converted to years as
\texttt{age\_at\_diagnosis / 365.25}. FIGO stage is mapped to an explicit
clinical ordinal with sub-stages collapsed to the major stage
($\mathrm{I}\!\to\!1$, $\mathrm{II}\!\to\!2$, $\mathrm{III}\!\to\!3$,
$\mathrm{IV}\!\to\!4$); an unmapped stage raises an error rather than being
silently coded. The source table holds 547 TCGA-OV records. Dropping rows whose
analysis fields carry the GDC missing marker leaves $n = 425$, and every number
reported in the paper is computed on those 425 rows; the 547 figure describes the
extracted table before missing-value filtering and is not an analysis sample size.
The split sizes follow directly: $318/107$ for the Task~1 train/test split and
$212/106/107$ for the Task~2 model-training/calibration/evaluation split.

Two design matrices are built from the same rows:

\begin{itemize}\tightlist
\item \textbf{Estimator design $X_{\mathrm{est}}$} ($p = 9$ columns, primary
four-field set): \texttt{age}, \texttt{stage\_ordinal}, and one-hot indicators
(\texttt{drop\_first=True}, lower-cased and stripped levels) for race
(\texttt{asian}, \texttt{black or african american},
\texttt{native hawaiian or other pacific islander}, \texttt{not reported},
\texttt{white}) and ethnicity (\texttt{not hispanic or latino},
\texttt{not reported}). Every method sees this identical design.
\item \textbf{Simulation design $X_{\mathrm{sim}}$} (4 columns):
$[\texttt{age},\ \texttt{stage\_ord},\ \texttt{race\_code},\
\texttt{ethnicity\_code}]$, with integer category codes for the two
categorical fields. Age stays in column 0 and stage in column 1 so the DGP's
$\delta(\cdot)$ and propensity remain functions of standardised age and stage.
\end{itemize}

\subsection*{A.2 Data-generating process}

$X_{\mathrm{sim}}$ is column-standardised,
$X_s = (X - \bar{X}) / (\sigma_X + 10^{-8})$, with $d = 4$. Coefficient vectors
are drawn \emph{once} from a fixed generator
(\texttt{numpy.random.RandomState(123)}), so they are identical across all
scenarios, seeds and repetitions:
\[
\beta_0 \sim \mathcal{N}(0, 0.5^2)^{d},\qquad
\beta_1 \sim \mathcal{N}(0, 0.5^2)^{d},\qquad
f = X_s\beta_0 .
\]
Let $a = X_{s,\cdot 0}$ (standardised age) and $s = X_{s,\cdot 1}$
(standardised stage), and $\sigma(z) = 1/(1+e^{-z})$. The four DGPs are:
\[
\begin{aligned}
\text{DGP 1 (constant): } &\ \delta = 0.5\cdot\mathbf{1}_n,\\
\text{DGP 2 (linear): } &\ \delta = 0.4\,a + 0.3\,s,\\
\text{DGP 3 (nonlinear): } &\ \delta = 1.0\,\sigma(X_s\beta_1) - 0.5,\\
\text{DGP 4 (nonlinear + interaction): } &\ \delta = 1.0\,\sigma(X_s\beta_1) - 0.5 + 0.4\,a\,s.
\end{aligned}
\]

\paragraph{Treatment assignment.}
\[
e(X) = \sigma(-0.6\,a - 0.4\,s),\qquad T_i \sim \mathrm{Bernoulli}(e_i),
\]
drawn as $T_i = \mathbf{1}\{u_i < e_i\}$ with $u_i$ from the per-run generator.
Coefficients $(-0.6, -0.4)$ and no intercept.

\paragraph{Potential outcomes and noise.}
Outcomes are simulated binary outcomes, interpreted in the narrative benchmark as five-year survival; they are not observed TCGA survival outcomes. There is no additive Gaussian noise term; all stochasticity is Bernoulli sampling:
\[
p_0 = \sigma(f),\quad p_1 = \sigma(f + \delta),\quad
Y_i(0)\sim\mathrm{Bernoulli}(p_{0i}),\quad Y_i(1)\sim\mathrm{Bernoulli}(p_{1i}),
\]
$Y_i = T_i Y_i(1) + (1-T_i)Y_i(0)$, and the estimand (ground truth for PEHE) is
$\tau^{\mathrm{true}} = p_1 - p_0$ on the probability scale. $Y(0)$, $Y(1)$ and
$T$ are drawn from a single \texttt{numpy.random.RandomState(seed)} in the order
$T$, $Y(0)$, $Y(1)$.

\subsection*{A.3 Random seeds and split logic}

\begin{itemize}\tightlist
\item Global seed in \texttt{code.py}: \texttt{SEED = 42} (\texttt{np.random.seed(42)}); default \texttt{LeafEncoder} seed.
\item DGP coefficient seed: \texttt{123} (fixed, shared across all runs).
\item \textbf{Task 1:} \texttt{seed} $\in \{0,\dots,9\}$ (10 seeds) $\times$ 4 DGPs $\times$ 9 methods $=$ 360 rows in \texttt{cate\_benchmark\_primary.csv}. The same \texttt{seed} value is used for the DGP draw, the split, and every estimator's \texttt{random\_state}, so all methods see an identical split.
\item \textbf{Task 1 split:} \texttt{train\_test\_split(test\_size=0.25, random\_state=seed, stratify=T)}, i.e.\ 75/25 train/test ($318$ train, $107$ test before trimming). Propensity trimming $0.05 < e_i < 0.95$ is applied \emph{to the training fold only}; the test fold is untrimmed.
\item \textbf{Task 2:} 30 repetitions, \texttt{rep} $\in \{0,\dots,29\}$, DGP scenario 4. Nested split, both stratified on $T$: first \texttt{train\_test\_split(test\_size=0.5, random\_state=rep)} $\rightarrow$ model-training ($212$) vs.\ rest ($213$); then \texttt{train\_test\_split(rest, test\_size=0.5, random\_state=rep+10000, stratify=T[rest])} $\rightarrow$ calibration $n=106$ / evaluation $n=107$. Effective ratio 50/25/25. Splitting happens \emph{before} any fitting; the $0.05$--$0.95$ trim again applies only to the model-training fold.
\item \textbf{Task 3:} state generator \texttt{numpy.random.default\_rng(7)}; expanded stress-test generator \texttt{random.Random(3)}; mock backend \texttt{random.Random(11)} with corruption probability $0.30$ (validation only, never reported as a live result).
\item Bootstrap seed: \texttt{numpy.random.default\_rng(0)}.
\end{itemize}

\subsection*{A.4 XGBoost settings (leaf-representation encoder)}

\texttt{xgboost.XGBClassifier} with \texttt{n\_estimators=100},
\texttt{max\_depth=4}, \texttt{learning\_rate=0.1},
\texttt{eval\_metric="logloss"}, \texttt{random\_state=seed},
\texttt{verbosity=0}; all other parameters at library defaults. It is fit on
$(X, Y)$ and used only via \texttt{.apply()} to extract per-tree leaf
assignments, which are one-hot encoded and concatenated to the raw covariates.

\paragraph{Cross-fitted encoder (variant C, primary).}
\texttt{StratifiedKFold(n\_splits=5, shuffle=True, random\_state=seed)} on the
training fold, stratified on $Y$ (or on a constant vector if $Y$ is degenerate).
For each fold an encoder is fit on the other $K-1$ folds and applied to the
held-out fold, so no held-out outcome enters its own representation. Leaf
columns are keyed by (tree index, exact split path) rather than by the numeric
leaf ID, because independent XGBoost fits do not share leaf numbering; the union
of keys across folds defines one shared column dictionary. For unseen data all
five fold-encoders are applied and their aligned one-hot matrices are averaged.
Variant B (\emph{leaky}) fits one encoder on the full training
fold; variant A uses raw covariates with no leaf features.

\subsection*{A.5 Base learners and estimator settings}

All nuisance and meta-learner base models are scikit-learn forests:
\begin{itemize}\tightlist
\item \texttt{make\_rf\_reg(seed)} $=$ \texttt{RandomForestRegressor(n\_estimators=200, max\_depth=6, min\_samples\_leaf=5, random\_state=seed)}.
\item \texttt{make\_rf\_clf(seed)} $=$ \texttt{RandomForestClassifier(n\_estimators=200, max\_depth=6, min\_samples\_leaf=5, random\_state=seed)}.
\end{itemize}
All other arguments are library defaults.

\begin{table}[H]
\centering\small
\caption{Estimator configurations (EconML 0.16.0 unless noted). Table~\ref{tab:main} reports the \texttt{NonParamDML} and \texttt{DRLearner} rows under these implementation names; \texttt{ForestDRLearner} is the doubly robust arm inside CaML-OP and is not a standalone baseline row.}
\begin{tabular}{@{}ll@{}}
\toprule
Method & Configuration \\
\midrule
S-Learner & \texttt{SLearner(overall\_model=make\_rf\_reg(seed))} \\
T-Learner & \texttt{TLearner(models=make\_rf\_reg(seed))} \\
X-Learner & \texttt{XLearner(models=make\_rf\_reg(seed),} \\
          & \quad \texttt{propensity\_model=make\_rf\_clf(seed))} \\
Causal Forest & \texttt{CausalForestDML(model\_y=make\_rf\_reg(seed),} \\
              & \quad \texttt{model\_t=make\_rf\_clf(seed), discrete\_treatment=True,} \\
              & \quad \texttt{n\_estimators=300, min\_samples\_leaf=5, cv=5,} \\
              & \quad \texttt{random\_state=seed)}; 95\% intervals via \\
              & \quad \texttt{effect\_interval(X, alpha=0.05)} \\
NonParamDML & \texttt{NonParamDML(model\_y=make\_rf\_reg(seed),} \\
(R-learner row) & \quad \texttt{model\_t=make\_rf\_clf(seed), model\_final=make\_rf\_reg(seed),} \\
          & \quad \texttt{discrete\_treatment=True, cv=5, random\_state=seed)} \\
DRLearner & \texttt{DRLearner(model\_propensity=make\_rf\_clf(seed),} \\
(DR-learner row) & \quad \texttt{model\_regression=make\_rf\_reg(seed),} \\
           & \quad \texttt{model\_final=make\_rf\_reg(seed), cv=5, random\_state=seed)} \\
ForestDRLearner & \texttt{ForestDRLearner(model\_regression=make\_rf\_reg(seed),} \\
(CaML-OP DR arm) & \quad \texttt{model\_propensity=make\_rf\_clf(seed), n\_estimators=300,} \\
                 & \quad \texttt{min\_samples\_leaf=5, cv=5, random\_state=seed)} \\
\bottomrule
\end{tabular}
\end{table}

\paragraph{Causal forest implementation.}
The causal forest is EconML's \texttt{CausalForestDML} (generalised random
forest / honest causal forest), not the R \texttt{grf} package; settings as
above.

\paragraph{CaML-OP ensemble.}
On the (possibly leaf-augmented) design, the R component
(\texttt{NonParamDML}) and DR component (\texttt{ForestDRLearner}) are fit on
the same training fold and combined with fixed equal weights,
$\hat{\tau} = 0.5\,\hat{\tau}_R + 0.5\,\hat{\tau}_{DR}$. The reported interval
is heuristic: the \texttt{ForestDRLearner} 95\% interval half-width is
multiplied by $1.10$ and re-centred on $\hat{\tau}$, i.e.
$\hat{\tau} \pm 1.10\,(\hat{h}_{DR}^{\mathrm{hi}} - \hat{h}_{DR}^{\mathrm{lo}})/2$.
The $1.10$ inflation is a stated convention acknowledging ensemble noise, not a
calibrated quantity.

\paragraph{Metrics.}
$\mathrm{PEHE} = \sqrt{\frac{1}{n}\sum_i(\tau_i - \hat{\tau}_i)^2}$;
$\mathrm{MAE} = \frac{1}{n}\sum_i |\tau_i - \hat{\tau}_i|$;
ATE error $= |\overline{\tau} - \overline{\hat{\tau}}|$. Sign accuracy is
$\frac{1}{n}\sum_i \mathbf{1}\{\operatorname{sign}\hat{\tau}_i = \operatorname{sign}\tau_i\}$,
reported as \texttt{NaN} (and excluded from means) whenever the true surface has
a single sign on that test set, which is why DGP 1 rows carry no sign
accuracy. Policy value uses the AIPW estimator with propensity clipped to
$[0.05, 0.95]$ and shared plug-in arm regressions $\hat{\mu}_0, \hat{\mu}_1$ fit
once per run on the training fold.

\subsection*{A.6 Bootstrap procedure}

Paired, non-parametric, over the $(\text{DGP}, \text{seed})$ pairing key.
For each contrast, per-pair PEHE differences $d_j$ are formed by aligning the
two methods on \texttt{(dgp, seed)}; $B = 10{,}000$ resamples of size $|d|$ are
drawn with replacement using \texttt{numpy.random.default\_rng(0)}; the
statistic is the resample mean; the interval is the percentile interval at
level $0.95$, i.e.\ the empirical $2.5\%$ and $97.5\%$ quantiles of the
bootstrap means. No bias correction or acceleration is applied. Contrasts
reported: CaML-OP[C\_crossfit] $-$ S-Learner, and CaML-OP[C\_crossfit] $-$
CaML-OP[A\_raw], with $n_{\text{pairs}} = 40$.

Repetition-level intervals in the conformal study are percentile intervals over
the 30 repetitions (empirical $2.5\%$/$97.5\%$ quantiles), not bootstrap
intervals.

\subsection*{A.7 Conformal calibration settings}

Split conformal on the held-out calibration fold, with the nonconformity score
formed against the \emph{known} semi-synthetic CATE (an oracle quantity,
available only because outcomes are simulated):
\[
s_i = \frac{|\tau_i - \hat{\tau}_i|}{\max(\hat{h}_i,\ \varepsilon)},
\qquad \varepsilon = 10^{-9},
\]
where $\hat{h}_i$ is the uncalibrated heuristic half-width from
Section~A.5. The threshold is the $k$-th order statistic of the calibration
scores,
\[
k = \min\!\big(\max(\lceil (n+1)\,(1-\alpha) \rceil,\ 1),\ n\big),
\qquad \hat{q} = s_{(k)},
\]
with the finite-sample $(n+1)$ correction. A plain
\texttt{np.quantile(scores, target)} is deliberately \emph{not} used. The
threshold is then frozen and the calibrated interval on the untouched
evaluation fold is $\hat{\tau}_i \pm \hat{q}\max(\hat{h}_i, \varepsilon)$.
Target levels: $0.80$, $0.90$, $0.95$, $0.99$. Abstention is defined as the
calibrated interval spanning zero. Calibration $n = 106$, evaluation
$n = 107$, 30 repetitions, DGP 4. Coverage is marginal only; no
subgroup-conditional guarantee is claimed.

\subsection*{A.8 Software environment}

Python 3.10.11 (Windows). Package versions as installed in the project virtual
environment at the time of the reported runs:

\begin{table}[H]
\centering\small
\caption{Package versions.}
\begin{tabular}{@{}ll@{}}
\toprule
Package & Version \\
\midrule
numpy & 2.2.6 \\
pandas & 2.3.3 \\
scikit-learn & 1.6.1 \\
scipy & 1.15.3 \\
statsmodels & 0.14.6 \\
xgboost & 3.2.0 \\
econml & 0.16.0 \\
lifelines & 0.30.0 \\
shap & 0.48.0 \\
matplotlib & 3.10.9 \\
seaborn & 0.13.2 \\
requests & 2.34.2 \\
\bottomrule
\end{tabular}
\end{table}

\subsection*{A.9 LLM experiment: models, date, decoding}

\begin{itemize}\tightlist
\item Exact model IDs: \texttt{claude-sonnet-5} and
\texttt{claude-haiku-4-5-20251001}, called through the Anthropic Messages API
(\texttt{https://api.anthropic.com/v1/messages}, header
\texttt{anthropic-version: 2023-06-01}).
\item Date the live experiments were run: \textbf{14 August 2026} (recorded in
the \texttt{run\_date} column of \texttt{llm\_verifier\_ablation.csv} for every
row of both models).
\item Requested temperature: $0.0$. \texttt{claude-haiku-4-5-20251001} accepted
it and its rows record \texttt{temperature = 0.0};
\texttt{claude-sonnet-5} rejected the \texttt{temperature} parameter, so the
harness dropped it and re-issued the call, and those rows record
\texttt{temperature = model\_default}. No other decoding parameters are set.
\item \texttt{max\_tokens}: 600 for typed (tool-call) conditions, 400 for the
free-text baseline. Request timeout 120\,s.
\item Evaluation set: 160 held-out synthetic patient states
(\texttt{make\_states(n=160, seed=7)}), stratified over the $3\times2\times2\times2$
grid \{positive, negative, near-zero effect\} $\times$ \{narrow, wide interval\}
$\times$ \{low, high R/DR disagreement\} $\times$ \{low, high robustness
warning\}. \texttt{must\_abstain} is defined deterministically as
$\mathrm{ci\_lo} \le 0 \le \mathrm{ci\_hi}$.
\item Conditions on identical states: \textbf{A} free-text baseline (no schema,
no gate), \textbf{B} typed tool-call output with no gate, \textbf{C} typed
output $+$ five deterministic checks (\texttt{max\_retries=0}), \textbf{D} typed
output $+$ five checks $+$ up to two targeted retries
(\texttt{max\_retries=2}). B, C and D share one typed generation for the first
attempt (\texttt{shared\_generation=1}) so that B vs.\ C isolates the gate
exactly; D's retries are fresh calls. Escalated outputs are never released, which
is what the \texttt{released\_unsupported} column records.
\item Free-text rubric validation: 30 \texttt{claude-sonnet-5} narratives
(\texttt{freetext\_validation.csv}).
\end{itemize}

\subsection*{A.10 System prompt (shortened exact version)}

The prompt is a single user message; there is no separate system prompt. The
template is reproduced below with the field-substitution placeholders
kept. A SHA-256 prefix of \texttt{PROMPT\_TEMPLATE + RETRY\_SUFFIX}
(\texttt{PROMPT\_VERSION}) is stamped into every result row so a stored pass
rate cannot silently go stale.

\begin{quote}\small\ttfamily\raggedright
You are a clinical decision-support narrative generator. Given the structured
outputs of a causal machine-learning model estimating the individualized effect
of platinum-based chemotherapy on five-year ovarian cancer survival, submit a
narrative via the submit\_clinical\_narrative tool.\\[4pt]
Rules (must follow exactly):\\
- "sign" must be "positive" if tau\_hat > 0 or "negative" if tau\_hat < 0
\{abstain\_rule\}\\
- "cited\_features" must contain only feature names drawn from the top drivers
listed below. Do not invent or include features not listed.\\
- "cited\_values" must tag every number you state in the narrative with which
quantity it represents, using exactly one of these tags: tau\_hat, ci\_lo,
ci\_hi, mu0, mu1, or shap:<feature> for a listed driver. Do not cite a number
under the wrong tag. You must cite both mu0 and mu1 [\dots]. If "sign" is
"positive" (benefit), mu1 must exceed mu0; if "negative" (harm), mu0 must
exceed mu1 [\dots].\\
- "dominant\_uncertainty\_source" must name whichever of the three uncertainty
signals below is largest for this patient specifically (not a generic answer):
"sampling", "model\_disagreement", or "identification" [\dots]. Mention this
dominant source in the narrative in plain language.\\
- The "narrative" field itself must be two to three sentences, must state the
sign/magnitude of the effect and its uncertainty interval in plain language,
must name the cited features, must mention the dominant uncertainty source, and
must not introduce outside medical knowledge or claims not derivable from the
numbers given.\\
\{abstain\_narrative\_rule\}\\[4pt]
Patient covariates: \{covariates\}\\
Estimated ITE (tau\_hat): \{tau\_hat:.3f\}\\
95\% uncertainty interval: [\{ci\_lo:.3f\}, \{ci\_hi:.3f\}]\\
Expected outcome if untreated (mu0): \{mu0:.3f\}\\
Expected outcome if treated (mu1): \{mu1:.3f\}\\
R-Learner component estimate (tau\_r): \{tau\_r:.3f\}\\
DR-Learner component estimate (tau\_dr): \{tau\_dr:.3f\}\\
Top features increasing benefit: \{top\_positive\}\\
Top features decreasing benefit: \{top\_negative\}
\end{quote}

\texttt{\{abstain\_rule\}} is inserted only when \texttt{must\_abstain} is
true: \emph{``-- EXCEPT: a calibrated reliability check on this estimate has
determined the sign is not statistically distinguishable from zero, so `sign'
must be `indeterminate' instead, regardless of the raw tau\_hat value above.''}
Otherwise it is a full stop. The matching
\texttt{\{abstain\_narrative\_rule\}} requires the narrative to say the effect
is not reliably distinguishable from no effect and to claim neither benefit nor
harm.

Condition A uses a deliberately different, unconstrained prompt: the same
opening sentence and the same numeric patient block, asking for ``a 2--3
sentence narrative for this patient in plain free text'', with no rule list, no
tool and no schema.

\subsection*{A.11 Typed output schema}

Anthropic tool \texttt{submit\_clinical\_narrative}, forced via
\texttt{tool\_choice = \{"type": "tool", "name": "submit\_clinical\_narrative"\}}:

\begin{quote}\small\ttfamily\raggedright
\{\\
\hspace*{1em}"narrative": string,\\
\hspace*{1em}"sign": string, enum ["positive", "negative", "indeterminate"],\\
\hspace*{1em}"cited\_features": array of string,\\
\hspace*{1em}"cited\_values": array of \{ "quantity": string, "value": number \},\\
\hspace*{2em}\quad required: ["quantity", "value"],\\
\hspace*{1em}"dominant\_uncertainty\_source": string, enum ["sampling",\\
\hspace*{2em}"model\_disagreement", "identification"]\\
\}\\
required: ["narrative", "sign", "cited\_features", "cited\_values",\\
\hspace*{2em}"dominant\_uncertainty\_source"]
\end{quote}

Legal \texttt{quantity} tags are exactly
\texttt{tau\_hat}, \texttt{ci\_lo}, \texttt{ci\_hi}, \texttt{mu0}, \texttt{mu1},
and \texttt{shap:<feature>} for each listed top driver. A response with no tool
call is recorded as \texttt{sign="invalid"} and fails the gate.

\subsection*{A.12 The five verifier checks}

Applied to the structured fields, never to prose
(\texttt{default\_verifiers(tol=0.15, atol=0.005)}):

\begin{enumerate}\tightlist
\item \textbf{sign.} If \texttt{must\_abstain}, requires
\texttt{sign == "indeterminate"}. Otherwise requires \texttt{"positive"} if
$\hat{\tau} > 0$, else \texttt{"negative"}.
\item \textbf{features.} Set membership: every cited feature must belong to the
declared top-driver set (union of \texttt{top\_positive} and
\texttt{top\_negative}); at least one feature must be cited; if the driver set
is empty, no feature may be cited. No text search is performed.
\item \textbf{numeric.} Every cited value must match the specific quantity its
tag claims, within a combined tolerance
$|v_{\text{cited}} - v_{\text{true}}| \le \max(0.15\,|v_{\text{true}}|,\ 0.005)$.
A tag that is not in the legal set fails immediately. The absolute floor
$0.005$ matches the three-decimal display precision used throughout the
pipeline, which a pure relative tolerance would be unforgivingly tight against
near $\hat{\tau} \approx 0$.
\item \textbf{counterfactual consistency.} Both \texttt{mu0} and \texttt{mu1}
must be cited. If \texttt{sign} is \texttt{"positive"}, the cited
$\mu_1 > \mu_0$; if \texttt{"negative"}, $\mu_1 < \mu_0$. Vacuously passes when
the claim is \texttt{"indeterminate"} or the state is \texttt{must\_abstain}.
\item \textbf{dominant uncertainty.} Must equal the deterministic argmax of
three normalised scores:
\[
\begin{aligned}
\text{sampling} &= \tfrac{\mathrm{ci\_hi}-\mathrm{ci\_lo}}{0.3}, \quad
\text{model\_disagreement} = \tfrac{|\tau_R - \tau_{DR}|}{0.05}, \\
\text{identification} &= \tfrac{2.0}{\max(\mathrm{E},\,10^{-6})},
\end{aligned}
\]
with $\mathrm{E}$ the VanderWeele--Ding E-value on the risk-ratio scale,
$\mathrm{E} = \mathrm{RR} + \sqrt{\mathrm{RR}(\mathrm{RR}-1)}$ where
$\mathrm{RR} = \mu_1/\mu_0$ inverted if below 1 and $\mu$ clipped away from
$\{0,1\}$. The thresholds $0.3$ (CI width), $0.05$ (component disagreement) and
$2.0$ (fragile E-value) are stated conventions, not empirically calibrated.
\end{enumerate}

\subsection*{A.13 Retry rule}

\texttt{generate\_and\_check} runs at most \texttt{max\_retries + 1} attempts
(condition C: \texttt{max\_retries=0}, one attempt; condition D:
\texttt{max\_retries=2}, up to three attempts). After a failing attempt the
feedback strings of \emph{only the failing} verifiers are concatenated and
appended to the prompt together with the rejected submission
(\texttt{asdict(previous)}) under the suffix: \emph{``Your previous submission
was rejected by an automated faithfulness check for this reason: \{feedback\}
Previous submission: \{previous\} Submit a corrected version via
submit\_clinical\_narrative that fixes this while still following all rules
above.''} The loop returns \texttt{NarrativeStatus.PASSED} on the first attempt
where all five checks pass, otherwise
\texttt{NarrativeStatus.ESCALATED}, a distinct state routed to human review,
never auto-displayed and never counted as released output.

\subsection*{A.14 Free-text parser and rubric (condition A)}

Condition A has no schema, so the narrative is labelled by a regular-expression
rubric (\texttt{freetext\_scoring.py}), evaluated in the order a human reader
would apply it. The rubric replaced an earlier bare-substring parser that
scored 29/30 validation narratives as \texttt{assert\_pos}, including four
unambiguous negative assertions.

\begin{itemize}\tightlist
\item \textbf{no\_direction}: the text states a point estimate but explicitly
flags that the interval includes zero or the result is inconclusive, and makes
no treatment recommendation. Hedge patterns include \emph{spans/crosses/includes
zero}, \emph{cannot (confidently) distinguish / rule out / determine / exclude},
\emph{inconclusive}, \emph{not statistically significant / definitive /
conclusive}, \emph{no meaningful benefit / effect / difference}, \emph{little to
no}, \emph{indeterminate}, \emph{low confidence}, \emph{substantial or
meaningful uncertainty}, \emph{could range from}, \emph{disagree on the
direction}.
\item \textbf{assert\_pos}: the text endorses treatment (\emph{strong /
favorable / reasonable candidate}, \emph{supports treatment},
\emph{recommendation favoring}, \emph{reinforcing the recommendation}) or states
a directional positive conclusion (\emph{meaningfully improves}, \emph{increases
(predicted) survival}, \emph{clinically meaningful benefit}, \emph{absolute
benefit of}, \emph{percentage-point increase}) without flagging that the
interval spans zero.
\item \textbf{assert\_neg}: the same, in the direction of harm (\emph{lower
predicted}, \emph{decrease in estimated survival}, \emph{worse expected},
\emph{reduction in five-year survival}, \emph{unfavorable predicted},
\emph{negative expected benefit}).
\end{itemize}

Precedence: an explicit endorsement overrides hedging, because recommending
treatment is a directional claim however many caveats surround it; otherwise
hedged $\rightarrow$ \texttt{no\_direction}; otherwise negative patterns before
positive; default \texttt{no\_direction}. Under the abstention rule a violation
is any \texttt{assert\_*} on a \texttt{must\_abstain} state; reporting a point
estimate with an explicit ``this interval spans zero'' caveat is \emph{not}
counted (the lenient reading). A stricter reading, in which any signed
directional statement counts as an assertion regardless of caveats
(\texttt{score\_strict}), is reported alongside it; bare negative numbers are
never used as evidence, since interval bounds such as ``$-19.8\%$ to
$+39.5\%$'' would otherwise read as negative claims.

\subsection*{A.15 Verifier stress tests}

Two suites, both deterministic and API-free. (i) The original 15-case
regression suite from \texttt{narrative\_research\_eval.py}. (ii) An
auto-generated suite of 128 cases: 14 corrupted cases for each of the eight
implemented failure classes (\texttt{wrong\_sign},
\texttt{abstention\_violation}, \texttt{hallucinated\_feature},
\texttt{wrong\_numeric\_value}, \texttt{unknown\_numeric\_tag},
\texttt{missing\_mu\_citation}, \texttt{arm\_ordering\_contradiction},
\texttt{wrong\_uncertainty\_source}) $=112$ cases that must be detected, plus 16
uncorrupted faithful cases that must pass. Corruptions are injected into the
deterministic \texttt{mock\_llm} output with \texttt{random.Random(3)}; classes
that are logically inapplicable to a state are excluded (e.g.\
\texttt{wrong\_sign} on an abstain state, \texttt{abstention\_violation} on a
non-abstain state).